\section{The primitive polytope and its interior maximum}\label{st:variation}
For a fixed strict expression let $P$ be the compact polytope
obtained by replacing all strict leaf and product inequalities
by their weak versions. Turns retain $y=1-x$, and closure
retains the complementary relation. Its relative interior is
the strict state set: a strict state exists, so none of the
listed proper inequalities is identically an equality.

\subsection{Integer coordinates and endpoints}
\begin{lemma}[Primitive coordinates]\label{st:lattice}
A closed tree with $H$ leaves has $K=2H-3$ free integer
coordinates after the forced closing product is removed.
Every remaining finite label has the form
\begin{equation}\label{st:offsets}
 r_w=NL_w(k/N)+\delta_w,\qquad k\in\ZZ^K,
\end{equation}
where $L_w$ is integral affine and $\delta_w$ is a fixed
integer. There is no residual congruence restriction.
Every strict real state is approached by admissible states
at all sufficiently large levels of both parities.
For an open tree the full dimension is $2H-1$ and its
fixed-output slice has dimension $2H-2$.
\end{lemma}
\begin{proof}
Keep a coordinate for every leaf output and every product
output. A unary turn simply replaces its input by $d-r$.
The closed tree has $H-2$ retained products, in addition
to its $H$ leaves. Its two root subtrees are disjoint,
so the closing relation can be solved for an output
coordinate in one subtree with coefficient $1$ or $-1$.
No identification of the two root coordinates has been
made. In particular, this procedure never creates $2r=d$.
This proves the dimension and primitivity assertions.
Replacing $d$ by $N-1$ gives \eqref{st:offsets}.

Round $d$ times the free coordinates of a strict state,
and reconstruct the eliminated coordinate exactly.
The errors in all labels are bounded for the fixed
expression. Every proper inequality has a positive
slack times $d$, so the rounded state is admissible
for all sufficiently large $N$. The strict leaf slack
also puts its labels below $M$, at either parity.
For a prescribed open output $r_N/d\to b$ in the actual
strict interval, choose strict states continuously
near $b$, fix that output first and round the other
coordinates. Such a continuous choice is obtained by
interpolating two nearby strict states with output
on either side of $b$.
\end{proof}

Define
\begin{align}
 J(t)&=\int_0^t\log(2\sin\pi u)\dd u,\qquad 0\le t\le1,
       \label{st:J}\\
 g(a)&=J(2a)-4J(a)
       =-2J(a)-2J(1/2-a),\label{st:g}\\
 m(x,y,z)&=J(x+y)+J(z)-J(x)-J(y)\notag\\
 &\hspace{11mm}-J(z-x)-J(z-y)-J(x+y-z).
       \label{st:m}
\end{align}
Here $J$ has its continuous endpoint values,
$J(1-t)=-J(t)$ and $J(1/2)=0$.
The complete action is the sum of one $g$ for each
original leaf and one $m$ for each retained product:
\begin{equation}\label{st:action}
 \A(y)=\sum_{\ell=1}^H g(a_\ell(y))
             +\sum_{v=1}^{H-2}m(x_v(y),y_v(y),z_v(y)).
\end{equation}
For an open expression the sum contains all its products.

The uniform factorial estimate
\cite[Lemma 5.2]{Newton} is
$\log S_r=NJ(r/N)+O(\log N)$ for $0\le r<N$.
Applied to \eqref{st:product} and \eqref{st:leaf}, it gives
\begin{equation}\label{st:weight-rate}
 \log w_N(\mathbf r)=N\A(\mathbf r/d)+O_E(\log N)
\end{equation}
uniformly on all nonzero positive states of a fixed
expression. In using this estimate at the endpoints,
the modulus of continuity of $J$ is
$O(\eta|\log\eta|)$. Thus the changes from $r/N$ to
$r/d$, and from $M/N$ to $1/2$, contribute
$O_E(\log N)$, even at the support boundary.

\subsection{The atoms and their boundary derivative}
For $a+b+c=1$ set
\[
 V(a,b,c)=-\pi\{J(a)+J(b)+J(c)\}.
\]
This is a nonnegative function, zero on the boundary of
the angle simplex and positive in its interior.
Only these elementary properties and its derivatives
are used here. With $c=1-a-b$ its Hessian has entries
\[
 -\pi^2
 \begin{pmatrix}
  \cot\pi a+\cot\pi c&\cot\pi c\\
  \cot\pi c&\cot\pi b+\cot\pi c
 \end{pmatrix}.
\]
Both displayed diagonal sums are positive, and the
determinant is $\pi^4$, since
$\cot\pi a\cot\pi b+\cot\pi b\cot\pi c+
\cot\pi c\cot\pi a=1$. The Hessian is therefore negative
definite at a positive triple.

At a product put
\[
 p=z-x,\quad q_0=z-y,\quad r=x+y-z,\quad w=1-x-y.
\]
The subscript on $q_0$ distinguishes this angle from
the root of unity. Direct substitution gives
\begin{align}
 \pi g(a)&=2V(a,1/2-a,1/2),\label{st:leaf-atoms}\\
 \pi m(x,y,z)&=
 V(y,x,w)+V(p,q_0,2r+w)+V(p+q_0,r,r+w).
 \label{st:gate-atoms}
\end{align}
All three entries of each atom are nonnegative and sum
to one on the original gate support. Consequently the
whole action is concave on its shared-label polytope.

There are two different kinds of flat atom. At
$(0,a,1-a)$ with $0<a<1$, the directional derivative
toward a positive atom is $+\infty$:
$-J(t)=-t\log t+O(t)$ near zero.
At $(0,0,1)$ the singularities cancel. More explicitly,
along the segment with small entries $t\alpha,t\beta$,
\[
 V(t\alpha,t\beta,1-t(\alpha+\beta))
 =\pi t\{(\alpha+\beta)\log(\alpha+\beta)
                  -\alpha\log\alpha-\beta\log\beta\}
       +O(t^2).
\]
The derivative is finite. A positive atom also has a
finite derivative. Thus no atom contributes $-\infty$
along a segment toward a strict state.

\begin{lemma}[Endpoint descent]\label{st:barrier}
At every nonstrict point of a nonempty strict closed
polytope, or of an open fixed-output slice with output
in its strict interval, at least one atom in
\eqref{st:leaf-atoms}--\eqref{st:gate-atoms} is flat
but is not a vertex of its angle simplex.
\end{lemma}
\begin{proof}
First suppose a proper wire has value zero or one.
A turn replaces it by its complementary endpoint.
At a product with output zero both inputs are zero.
At a product with output one the inputs sum to one:
if both are interior, the first atom is $(y,x,0)$,
as required; otherwise one may descend to an endpoint
child. Repeating this descent in the finite tree
terminates at a leaf endpoint. A leaf of value zero
has atom $(0,1/2,1/2)$; a leaf cannot have value one.
The same argument applies to either child of the
deleted root: one descends inside that child, without
putting the removed flat gate back.

Suppose now that every wire is in $(0,1)$.
A leaf at its allowed upper endpoint $1/2$ gives
$(1/2,0,1/2)$. If every leaf is strict, a nonstrict
state must lie on a retained gate face.
When $w=0$, the first atom $(y,x,0)$ is nonvertex-flat.
When $r=0$, the third atom is $(z,0,1-z)$.
When $p=0$ and $q_0>0$, the second atom has one zero
entry and two positive entries; the case $q_0=0<p$
is symmetric. In the remaining intersection
$p=q_0=0$, one has $x=y=z\in(0,1)$ and the third
atom is $(0,x,1-x)$. This lists every gate boundary.

For a fixed-output slice the output itself is in
$(0,1)$, and the same descent applies to each proper
endpoint wire. The segment used below keeps that
output fixed. No boundary caused only by imposing
the output is included among these proper faces.
\end{proof}

\begin{theorem}[Unique strict optimizer]\label{st:positive-action}
Every strict closed expression has a unique maximizer
$y_*$ of its complete action. It is strict and its
Hessian in the primitive coordinates is negative
definite. The positive scalar satisfies the all-level
limit
\begin{equation}\label{st:positive-limit}
 \lim_{N\to\infty}\frac1N\log Z_N^0=F_0,\qquad
 F_0=\A(y_*)>0 .
\end{equation}
For an open expression of nonempty type, each output
$b$ in its strict interval has a unique strict
fixed-output maximizer. The profile
\[
 f_E(b)=\max_{\text{output }b}\A_E
\]
is real analytic and strictly concave on that interval.
It has a continuous extension to the closed feasible
output interval.
\end{theorem}
\begin{proof}
Compactness gives a maximizer. Join a nonstrict point
to a strict state in the same polytope. By
Lemma~\ref{st:barrier}, one atom contributes $+\infty$
to the one-sided derivative, and none contributes
$-\infty$. Hence the action increases and the point
cannot maximize.

At a strict product the gate Hessian is already
negative definite in $(x,y,z)$. Indeed the first
atom in \eqref{st:gate-atoms} detects
$(\delta x,\delta y)$; if those vanish, the second
detects $\delta z$. Its third atom is nonpositive
and cannot cancel either negative term.
The leaf Hessian is strictly negative since
\[
 g'(a)=2\log\cot\pi a,\qquad
 g''(a)=-\frac{4\pi}{\sin2\pi a}<0 .
\]
If all atom variations vanish, every leaf output
and retained product output is fixed. Turns then
fix their output wires, and the closing relation
fixes the eliminated coordinate. Thus the sum is
negative definite on the primitive tangent space.
This proves uniqueness by concavity.

For an open expression the same argument applies
including its output coordinate: its last product
detects that coordinate, or the expression is a
turned leaf and its leaf detects it. On a fixed-output
slice, the internal Hessian is therefore invertible.
The analytic implicit-function theorem gives a
real analytic maximizing state and profile.
The full negative Hessian has a negative Schur
complement in the output coordinate, so $f_E''<0$.
For the single-leaf case this is the displayed
formula and its reflection.

Continuity at a feasible endpoint is worth recording.
For any endpoint maximizing state, interpolate it
with an interior strict state. This supplies states
at every nearby interior output, and their actions
converge to the endpoint maximum. Conversely, every
sequence of maximizing states has a convergent
subsequence in the compact full polytope; its limit
is feasible at the endpoint. These two observations
prove continuity. They do not require finite support
at the exact endpoint for every parity.

There are at most $N^K$ positive states.
Equation~\eqref{st:weight-rate} gives the upper
limit in \eqref{st:positive-limit}. Primitive
rounding of $y_*$ gives a genuine positive summand
at every sufficiently large $N$ and the matching
lower limit. The strict leaf atoms are positive,
so $F_0>0$.
\end{proof}

\begin{lemma}[Gate curvature without leaf curvature]\label{st:gate-curvature}
For a closed strict tree with $H\ge3$, the sum of the
retained gate Hessians is itself negative definite
in the primitive coordinates.
\end{lemma}
\begin{proof}
Contract unary turn chains. Every wire in a nontrivial
root subtree appears at a retained product, either
as an input or as an output. If all gate variations
vanish, these wires have zero variation. The only
possible uncovered root subtree is a single leaf
with turns; its variation is fixed by the complementary
root relation and the other subtree. Both subtrees
can be single leaves only when $H=2$. The injectivity
and the strict gate Hessian proved above give the result.
\end{proof}
